%! Author = Sujal Singh
%! Date = 10/23/23

% Preamble
\documentclass[11pt,a4paper,oneside]{scrarticle}

% Packages
\usepackage{mathptmx}
\usepackage{tabularray}
\usepackage{graphicx}
\usepackage{sectsty}
\usepackage{titlesec}
\usepackage{setspace}
\usepackage{romannum}
\usepackage{hyperref}

\renewcommand*{\maketitle}{
    \pagenumbering{gobble}
    \begin{center}
        \vspace*{8em}
        \includegraphics[width=200pt]{../../tools/pdf/template/logo} \\[40pt]
        \textbf{\Huge Law Assignment} \\[40pt]
        \begin{table}[htb]
            \label{tab:name-slip}
            \begin{tblr}{rows={30pt},colspec={|X|X[3]|},hlines}
                \textbf{Course}            & Industrial Internet of Things \\
                \textbf{Batch}             & 1                             \\
                \textbf{Name}              & Sujal Singh                   \\
                \textbf{Enrollment Number} & XXXXXXXXXXXX                  \\
            \end{tblr}
        \end{table}
    \end{center}
    \newpage
    \pagenumbering{arabic}
}

\linespread{1.5}
\titleformat{\section}{\Large\bfseries\centering}{\thesection.}{0.5em}{}

% Document
\begin{document}
    \maketitle

    \begin{center}
        \textbf{\LARGE The Indian Government's Stance on Fundamental Rights in the Digital World.}
        \\[10pt]
    \end{center}
    \begin{flushright}
        \textit{--- Sujal Singh} \\[20pt]
    \end{flushright}

    %------------------------------------------------------------------------------------------------------------------%
    %------------------------------------------------------------------------------------------------------------------%

    \section*{Abstract}\label{sec:abstract}
    The Internet is one of mankind's greatest achievements and an open Internet, free of censorship and control by any
    authority, is crucial to avoid a dystopian future, but ever since it's inception, government authorities have
    employed various tactics to censor the information published via the World Wide Web.\ This report details the
    state of fundamental rights as they exist in the digital landscape of India and how the Indian government has often
    violated these rights to censor the Internet and use it as a tool for mass-surveillance.\ Such acts go against
    constitutional values as envisioned by the authors of the Indian Constitution, and people in relevant fields of
    studies should feel responsible to oppose the continued decline of the open Internet, and they should push for
    technical solutions that uphold constitutional values.


    %------------------------------------------------------------------------------------------------------------------%
    %------------------------------------------------------------------------------------------------------------------%


    \section{Introduction}\label{sec:introduction}

    The Indian Constitution provides six fundamental rights to all citizens of India, one of which is the right to
    freedom.\ According to article 19(1)(a):
    \begin{quote}
        \singlespacing
        (1)
        All citizens shall have the right---\\
        \vspace*{-15pt}
        \begin{enumerate}
        (a)
            to freedom of speech and expression;\\
            \ldots
        \end{enumerate}
    \end{quote}
    citizens of India are provided a fundamental right to express their views relating to any matter.\ Due to this
    fundamental right, citizens are allowed to offer criticism to the state, in fact, without criticism the government
    loses accountability, and yet, various articles offering such criticism published on the web often has the
    possibility of landing the author in jail.
    \\
    According to the Supreme Court of India's ruling in the
    \textit{Justice K. S. Puttaswamy (Retd.) \& Another. vs. Union Of India \& Others} case, the right to privacy is
    protected as an intrinsic part of the right to life and personal liberty under Article 21 and as a part of the
    freedoms guaranteed by Part \Romannum{3} of the Constitution.
    \\
    These fundamental rights are guaranteed for every citizen of India, but in reality, the Indian government itself has
    a history of openly violating these rights.\ The following sections describe how the Indian government tackles
    various issues relating to maintaining fundament rights on the Internet.
\end{document}